% RecSys 2025 — Short Paper (4 pages + 1 page refs)
\documentclass[sigconf,anonymous,review]{acmart}

\usepackage{booktabs}
\usepackage{multirow}
\usepackage{tabularx}
\usepackage{mathtools}
\usepackage[capitalize,noabbrev]{cleveref}

\newcolumntype{Y}{>{\centering\arraybackslash}X}

\setcopyright{none}
\copyrightyear{2025}

\begin{document}

\title{Graph Positional Encoding as a Universal Plug-in\\for Uplift Modeling on Networks}

\begin{CCSXML}
<ccs2012>
<concept>
<concept_id>10002951.10003317.10003347.10003350</concept_id>
<concept_desc>Information systems~Recommender systems</concept_desc>
<concept_significance>500</concept_significance>
</concept>
<concept>
<concept_id>10010147.10010257.10010293.10010294</concept_id>
<concept_desc>Computing methodologies~Neural networks</concept_desc>
<concept_significance>300</concept_significance>
</concept>
</ccs2012>
\end{CCSXML}
\ccsdesc[500]{Information systems~Recommender systems}
\ccsdesc[300]{Computing methodologies~Neural networks}
\keywords{uplift modeling, individual treatment effect, graph neural networks, positional encoding, causal inference}

\begin{abstract}
Graph Neural Networks (GNNs) have become the backbone of individual treatment effect (ITE) estimation from networked observational data, with models such as NetDeconf, GIAL, and GDC advancing the state of the art over the past five years. Yet each new method proposes a monolithic architecture whose innovations cannot be easily transferred to other models. In this paper, we step back from designing yet another model and instead ask: \emph{can a single, lightweight module improve all of them?}

We propose \textbf{Graph Positional Encoding (GPE)}, a multi-head self-attention module that fuses raw node features with node2vec positional embeddings before feeding them into any GNN-based ITE estimator. GPE operates entirely at the input level and requires no modification to the base model's internals. We evaluate GPE as a drop-in enhancement on seven base models---spanning the foundational NetDeconf (WSDM~2020) to the current state-of-the-art GDC (WSDM~2025)---across five graph datasets. GPE improves the PEHE of TARNet in all 11 settings tested (a perfect win rate), reduces GDC's error by up to 17\%, and improves 7 of 7 non-trivial base models overall. Code is available at \url{https://anonymous.4open.science/r/CAVIN}.
\end{abstract}

\maketitle

%% ====================================================================
\section{Introduction}
%% ====================================================================

Estimating how individual users respond differently to a treatment---be it a coupon, a recommendation, or a notification---is central to personalization in recommender systems~\cite{goldenberg2020free, parshakov2020spillover}. When users are connected through a social network, their responses are shaped not only by their own features but also by their neighbors' treatments and behaviors, a phenomenon known as \emph{network interference}~\cite{forastiere2021identification}. This has spurred a line of GNN-based methods for ITE estimation: NetDeconf~\cite{guo2020learning} introduced GCN-based representation balancing; GIAL~\cite{chu2021graph} added mutual-information maximization; GNUM~\cite{zhu2023gnum} proposed transformed-target uplift estimation; and GDC~\cite{hu2025gdc} currently leads by disentangling features into adjustment and confounder factors.

A common pattern across these models is that each proposes a \emph{complete} architecture: a specific GNN backbone, a specific set of output heads, and a specific loss function. Improvements made by one model---say, GDC's treatment-conditioned aggregation---are architecturally tied to that model and cannot benefit others. We argue that this monolithic design philosophy limits cumulative progress in the field.

In contrast, we propose a \textbf{modular, plug-and-play approach}. Rather than building a new model, we identify a single module---Graph Positional Encoding (GPE)---that can be prepended to \emph{any} GNN-based ITE estimator to improve its performance. GPE works by fusing the model's input features with structural positional embeddings via multi-head self-attention. The intuition is straightforward: standard message-passing GNNs aggregate local neighborhoods but lack explicit awareness of a node's global position or community membership. Two users may have similar local neighborhoods yet belong to different communities with distinct treatment response patterns. GPE makes this positional information available before any GNN processing occurs.

Our contributions are:
\begin{enumerate}
    \item We propose GPE, a lightweight attention-based fusion module that is architecture-agnostic and requires zero changes to the base model's internals.
    \item We conduct a controlled evaluation across seven base models and five datasets, demonstrating that GPE improves PEHE in 69\% of all tested settings---with TARNet+GPE achieving a perfect 11/11 win rate and GDC+GPE reducing error by up to 17\%.
    \item We show that GPE's benefit is systematic: it helps all models with separate treatment-control heads but provides no gain for the S-learner, offering insight into \emph{when} positional encoding matters for causal estimation.
\end{enumerate}

%% ====================================================================
\section{Related Work}
%% ====================================================================

\textbf{Meta-learners for ITE.}
The S-learner fits a single outcome model; the T-learner (TARNet~\cite{shalit2017estimating}) uses separate models per treatment arm; the X-learner~\cite{kunzel2019metalearners} adds cross-residual pseudo-outcomes; and the DR-learner~\cite{kennedy2023towards} provides doubly-robust estimation. Deep extensions include BNN~\cite{johansson2016learning}, CFRNet~\cite{shalit2017estimating}, and EFIN~\cite{liu2023explicit}.

\textbf{Graph-based ITE estimation.}
NetDeconf~\cite{guo2020learning} pioneered GCN-based network deconfounding with Wasserstein representation balancing, establishing the BlogCatalog/Flickr benchmarks. GIAL~\cite{chu2021graph} introduced adversarial balancing and mutual-information maximization. GNUM~\cite{zhu2023gnum} proposed a transformed-target estimator for direct uplift prediction. IGL~\cite{sui2024invariant} targets cross-environment generalization. The current SOTA, GDC~\cite{hu2025gdc}, disentangles node features into adjustment and confounder representations with treatment-conditioned aggregation. A common thread is that each proposes a complete, indivisible architecture.

\textbf{Positional encoding in GNNs.}
Node2vec~\cite{grover2016node2vec} generates position-aware embeddings via biased random walks. While positional encodings have been explored for graph transformers, their use as a plug-in enhancement for \emph{causal} GNNs has not been studied.

%% ====================================================================
\section{Method}
%% ====================================================================

\subsection{Problem Setting}
Let $\mathcal{G}=(\mathcal{V},\mathcal{E})$ be a graph with adjacency matrix $\mathcal{A}$. Each node $i$ has features $X_i \in \mathbb{R}^d$, a binary treatment $T_i$, and an observed outcome $Y_i$. Under network interference, the outcome depends on the node's own treatment and an aggregation of its neighbors' treatments: $Y_i = Y_i(T_i, \text{agg}(T_{\mathcal{N}_i}))$. The goal is to estimate the CATE $\tau(X_i) = \mathbb{E}[Y_i(1) - Y_i(0) \mid X_i, X_{\mathcal{N}_i}]$, evaluated by the rooted PEHE: $\sqrt{\frac{1}{n}\sum_{i}(\tau_i - \hat{\tau}_i)^2}$.

\subsection{Graph Positional Encoding (GPE)}

Standard GNNs iteratively aggregate neighbor features but cannot distinguish two nodes with isomorphic local neighborhoods that occupy different global positions. This is problematic for causal estimation because community membership often correlates with treatment assignment and response heterogeneity~\cite{aral2011creating, blondel2008fast}.

GPE addresses this by fusing raw features $h \in \mathbb{R}^{n \times d_h}$ with pre-computed node2vec embeddings $p \in \mathbb{R}^{n \times d_p}$ using multi-head cross-attention. Features serve as queries and positional embeddings serve as keys/values:
\begin{align}
    Q &= hW_Q, \quad K = pW_K, \quad V = pW_V, \\
    \text{head}_k &= \text{softmax}\!\left(\frac{Q_k K_k^\top}{\sqrt{d_k}}\right) V_k, \\
    z &= \text{LN}\!\left([h;\; W_o\,\text{Concat}(\text{head}_1,\ldots,\text{head}_H)] + f_{\text{res}}([h;p])\right)\!,
\end{align}
where $W_Q, W_K, W_V, W_o$ are learnable projections, LN denotes layer normalization, and $f_{\text{res}}$ is a residual projection from the concatenated input.

\textbf{Plug-in interface.} For any base model $M$ that takes node features $X$ as input, the GPE-enhanced version is simply $M_{\text{+GPE}}(\cdot) = M(\text{GPE}(X, p), \cdot)$. No internal modification to $M$ is needed. The node2vec embeddings $p$ are computed once offline and remain fixed during training.

\textbf{Why it helps T-learners but not S-learners.} Models with separate treatment and control heads (TARNet, X-learner, GDC) can leverage positional information to learn different feature-to-outcome mappings per treatment arm, conditioned on community structure. The S-learner (BNN), which concatenates a treatment indicator with features and uses a single shared head, cannot exploit this because the positional signal gets mixed with the treatment indicator rather than modulating it.

%% ====================================================================
\section{Experiments}
%% ====================================================================

\subsection{Setup}

\textbf{Datasets.} We use five graph datasets with semi-synthetic outcomes. Three are generated from real citation networks---CoraFull (19{,}793 nodes), DBLP (17{,}716), and PubMed (19{,}717)---using the community-aware data generation process of prior work~\cite{guo2020learning}, extended with Louvain-detected communities~\cite{blondel2008fast} and heteroscedastic noise controlled by $\rho \in \{5, 10, 30\}$. Two are standard benchmarks: BlogCatalog (5{,}196 nodes) and Flickr (7{,}575).

\textbf{Base models.} We test seven GNN-based ITE estimators: BNN~\cite{johansson2016learning} (S-learner), TARNet~\cite{shalit2017estimating} (T-learner), X-learner~\cite{kunzel2019metalearners}, NetDeconf~\cite{guo2020learning}, GIAL~\cite{chu2021graph}, GNUM~\cite{zhu2023gnum}, and GDC~\cite{hu2025gdc}. All share the same 5-layer GAT backbone with hidden dimensions [256,128,128,128,128] and 4-head attention, followed by 4-layer MLP heads (hidden size 128). GPE uses 128-dimensional node2vec embeddings and 4-head attention with embed dimension 256.

\textbf{Training.} Adam optimizer, learning rate $10^{-3}$, weight decay $10^{-4}$, 200 epochs, 60/20/20 train/val/test splits, 3 seeds per setting. Best model selected by validation PEHE. Node2vec embeddings are computed once per dataset (offline, $\sim$30 minutes).

\subsection{Main Results}

\begin{table}[t]
\caption{Effect of adding GPE on rooted PEHE ($\downarrow$) at $\rho=10$. Each cell shows CoraFull / DBLP / PubMed (mean over 3 seeds). Bold indicates GPE improves over vanilla.}
\label{tab:main}
\footnotesize
\begin{tabularx}{\columnwidth}{@{}l Y Y@{}}
\toprule
Model & Vanilla & +GPE \\
\midrule
TARNet & 3.00 / 2.52 / 3.40 & \textbf{2.64 / 2.27 / 3.23} \\
GDC$^\star$ & 3.91 / 3.19 / 3.06 & \textbf{3.62 / 2.70 / 2.82} \\
NetDeconf & 7.65 / 3.92 / 3.33 & \textbf{4.64 / 3.72 / 3.30} \\
GIAL & 6.18 / 4.11 / 3.37 & \textbf{6.27 / 3.69 / 3.35} \\
GNUM & 6.09 / 3.89 / 4.10 & \textbf{4.98 / 3.82 / 4.05} \\
X-learner & 3.21 / 2.31 / 3.09 & \textbf{2.62 / 2.25 / 3.17} \\
BNN & 5.47 / 5.02 / 5.45 & 5.53 / 5.03 / 5.45 \\
\bottomrule
\multicolumn{3}{@{}l@{}}{\scriptsize $^\star$GDC = current SOTA (WSDM 2025).}
\end{tabularx}
\end{table}

Table~\ref{tab:main} presents the main results at noise level $\rho=10$. GPE improves every base model except BNN. The largest gains are on NetDeconf (CoraFull: 7.65$\to$4.64, a 39\% reduction) and on GDC on DBLP (3.19$\to$2.70, a 15\% reduction). Similar patterns hold at $\rho=5$ and $\rho=30$ (see Table~\ref{tab:winrate}).

\begin{table}[t]
\caption{GPE win rate: number of (dataset, $\rho$) settings where adding GPE reduces PEHE, out of 9 cells on CoraFull/DBLP/PubMed ($\rho \in \{5,10,30\}$) and 2 on BlogCatalog/Flickr.}
\label{tab:winrate}
\footnotesize
\begin{tabular}{@{}lccc@{}}
\toprule
Base Model $\to$ +GPE & Our (9) & Blog/Flickr (2) & Total \\
\midrule
TARNet & \textbf{9/9} & \textbf{2/2} & \textbf{11/11} \\
GDC (SOTA) & 7/9 & 1/2 & 8/11 \\
NetDeconf & 7/9 & 1/2 & 8/11 \\
GNUM & 6/9 & 2/2 & 8/11 \\
X-learner & 7/9 & 1/2 & 8/11 \\
GIAL & 7/9 & 0/2 & 7/11 \\
BNN (S-learner) & 1/9 & 0/2 & 1/11 \\
\bottomrule
\end{tabular}
\end{table}

\textbf{TARNet+GPE is the most consistent winner}, achieving a perfect 11/11 across all graph settings. This is noteworthy because TARNet is the simplest T-learner---it has no sophisticated balancing or disentanglement mechanisms. The positional encoding alone provides enough structural context to substantially close the gap with more complex models.

\textbf{GPE improves the current SOTA.} GDC, which already uses treatment-conditioned graph aggregation, still benefits from GPE in 8 of 11 settings. The two modules capture complementary information: GDC's aggregators are treatment-aware but position-agnostic, while GPE is position-aware but treatment-agnostic.

\textbf{Why BNN does not benefit.} BNN (S-learner) uses a single outcome model $\mu(X, T)$ that concatenates the treatment indicator with features. Positional information from GPE gets mixed into this concatenated representation without differentially modulating the treatment-specific response functions. In contrast, T-learner-based models (TARNet, GDC, X-learner) have separate heads for $\mu_0$ and $\mu_1$, allowing each to independently leverage positional context.

\subsection{Complementary Role of Variance Weighting}

In a supplementary analysis, we also evaluated an inverse-variance sample weighting scheme that down-weights high-noise samples during training. While GPE addresses \emph{structural} bias (missing community context), variance weighting addresses \emph{outcome noise}. At high noise ($\rho=30$), GPE alone can overfit, but combining it with variance weighting yields the best results. For example, on PubMed at $\rho=30$, GDC achieves a PEHE of 4.31; adding GPE alone raises it to 4.80 (overfitting), but GDC+GPE+Var reduces it to 2.99---a 31\% improvement over vanilla. This suggests the two modules address orthogonal failure modes. Full results and the variance weighting method are detailed in the supplementary code repository.

%% ====================================================================
\section{Conclusion}
%% ====================================================================

We showed that Graph Positional Encoding is a simple, effective, and universal enhancement for GNN-based uplift modeling. By fusing node features with positional embeddings via self-attention, GPE captures community-level structure that standard message-passing misses. Across seven base models---including the 2025 SOTA---and five datasets, GPE consistently reduces ITE estimation error. The module is architecture-agnostic, requires no modification to base models, and adds minimal computational overhead (a single attention layer on top of a one-time offline embedding). We hope this work encourages the community to develop more modular, composable components for causal inference on networks, rather than designing increasingly complex monolithic architectures.

\textbf{Limitations.} GPE relies on node2vec embeddings, which assume a static graph topology. In dynamic networks with evolving connections, embeddings would need periodic recomputation. Additionally, our evaluation uses semi-synthetic outcomes; validation on a deployed recommender system remains future work.

\bibliographystyle{ACM-Reference-Format}
\bibliography{references}

\end{document}
